% Imporditud: tools/import_eksperimendid.py
% Allikas: Eesti füüsikaolümpiaadi eksperimendi ülesannete
%   kogu (Rael Kalda, 2023), ülesanne Ü20, lahendus L20.
\setAuthor{EFO žürii}
\setRound{piirkonnavoor}
\setYear{2012}
\setNumber{P E2}
\setDifficulty{1}
\setTopic{dünaamika}
\setVahendid{mutter, jupp niiti, mõõtejoonlaud, stopper, statiiv koos kinnitusklambriga}
\setPunktid{10}
\setAllikas{Eksperimendi kogu Ü20, lahendus lk 40}
\setLahendusseis{toores}

\prob{Pendel}
On teada, et niidist ja selle otsa kinnitatud väikesest kehast koosneva pendli
võnkeperiood $T$ ja selle niidi pikkus $l$ on väikese amplituudiga võnkumiste
korral seotud valemiga $T = \alpha\sqrt{l}$. Määrata võrdetegur $\alpha$
võimalikult täpselt.

\hint

\solu
Katse läbiviimiseks tuleb niidijupp siduda mutri külge ja niidi teine ots fikseerida

kinnitusklambri vahele.

Seejärel tuleb mõõta niidi pikkus (mutri massikeskmest kinnituskohani) ja mõõta

võnkeperiood. Võnkeperioodi täpseks määramiseks tuleb stopperiga mõõta aeg,

mis kulub mitme võnke tegemiseks (> 5), ja jagada see võngete arvuga.

Konstandi $\alpha$ võimalikult täpseks hindamiseks tuleb katseid korrata. (Niidi pikkust

võib, kuid ei pea varieerima.) teooriast $\alpha$ = $\pi$/g $\approx$ 2 s/$m_0$,5, korrektselt läbi Vastus tuleb matemaatilise pendli

viidud katse tulemus ei erine sellest rohkem kui 10\%.

Hindamine: Katse idee (2 p). Võnkeperioodi mõõtmisel n võnkeks (> 5) kulunud aja lugemine ja sellest võnkeperioodi arvutamine (1 p). Kui on mõõdetud ainult üheks võnkeks kulunud aeg, siis seda punkti ei saa. Korduskatsed (3x1 p) (iga korduskatse annab ühe punkti, maksimaalselt 3 ).

Arvutus $\alpha$ = T iga katseseeria jaoks (3x0,5 p), mitme mõõtmise korral $\alpha$ keskmis-

l

tamine (0,5 p).

Realistlik vastus ($\pm$10\%) (2 p).
\probend
